\documentclass[10pt]{letter}
\usepackage[margin=0.82in]{geometry}
\usepackage[hidelinks]{hyperref}
\signature{Nuonan Ouyang\\Corresponding Author}
\address{College of Science and Engineering\\James Cook University\\Australia\\nuonan.ouyang@my.jcu.edu.au}
\begin{document}
\begin{letter}{Professor Ahmed Louri\\Editor-in-Chief\\IEEE Transactions on Sustainable Computing}
\opening{Dear Professor Louri,}

Please find enclosed the major revision of manuscript TSUSC-2026-07-0200, ``Shielded Q-learning for Thermally-Safe, Energy-Aware Multi-Model Intrusion Detection Scheduling on Lightweight Edge Devices.''

We have revised the manuscript extensively in response to the Associate Editor and reviewers. The revision retains the submitted FSSQL-R method and operating constraints, but replaces the earlier replay-centered deployment evidence with new hardware experiments. The empirical implementation was rebuilt from the submitted method definitions and frozen before evaluation; the revised tables are replacement evidence rather than a claim of numerical reproduction of the earlier replay tables. The study now includes independent calibration on Raspberry Pi 3B+ and Raspberry Pi 4B 8GB, online closed-loop execution, a no-training Threshold controller, matched-seed statistics, Q-learning stability analysis, controller/switching overhead, direct POWER-Z KM003C DC input-energy measurement, and one-time fixed-threshold evaluation on the official UNSW-NB15 test set.

The new evidence led us to narrow several claims. We no longer present shielding as universally necessary, Q-learning as universally superior, or the guard as an unconditional complete-pipeline safety guarantee. We retain adverse results, including the stronger Safe-Greedy result on Pi 3B+, the strong Threshold baseline on Pi 4B, and the absence of a consistent physical-energy advantage for FSSQL-R. We also corrected the submitted protocol's use, as described, of test-set detector statistics in the scheduler objective: the revised reward uses validation-partition TPR/FPR only, and the official test set remained sealed until all development and hardware decisions were frozen. Recent 2024--2025 TSUSC, TMC, and TIFS work has also been added to strengthen the positioning.

A detailed point-by-point response is provided as a separate Summary of Changes document, and supplemental material contains additional calibration, statistical, energy, and official-test details.

Thank you for your consideration of the revised manuscript.

\closing{Sincerely,}
\end{letter}
\end{document}
